\chapter{A File of C Written as Zeno}\label{ch-convert-c}

\verb|zeno convert| reads C as well as Python. The promise is the one
of the last chapter: what it writes is never silently wrong. What it
can say in Zeno it says in Zeno, with the numbers C would give, bit
for bit, and what it cannot it keeps as C, in an \verb|inline c| block
that the Zeno calls.

\section{A first file}

Here is a small meter, in C:

\begin{zeno}[file=meter.c, dir=convertc, src=tools/convert/meter.c, lines=2-59, run=no]
\end{zeno}

One command writes \verb|meter.zeno| beside it:

\begin{zeno}[file=read.txt, dir=convertc, show=no, cmd=zeno convert meter.c]
meter.c
\end{zeno}

This is the Zeno it wrote:

\begin{zeno}[file=read.txt, dir=convertc, show=no, cmd=cat meter.zeno]
meter.zeno
\end{zeno}

\verb|tick| keeps a count from one call to the next in a
\verb|static| variable, which Zeno has no word for, so it stays C and
\verb|main| calls it as \verb|c_part.tick()|. The rest is Zeno. The
struct is a class and the arrays are lists. \verb|sum += bytes[i]| on
a \verb|uint8_t| goes round at 256 as it does in C, written
\verb|u8.wrap|. \verb|total / count| takes the quotient toward zero as
C does, through a short helper with a comment that says why. The two
pointers that \verb|limits| writes through are lists of one: the
caller puts \verb|low| in \verb|low_box|, makes the call, and takes it
out again.

Both print the same:

\begin{zeno}[file=read.txt, dir=convertc, show=no, cmd=zeno run meter.zeno ;; cc -o meter meter.c ;; ./meter]
meter.zeno
\end{zeno}

\section{The numbers are C's}

The converter reads each type as C does on a 64-bit machine, and
writes the type of Zeno that gives the same sums:

\begin{center}
\begin{tabular}{lP{7cm}}
\toprule
C & Zeno \\
\midrule
\verb|uint8_t| to \verb|uint64_t| & \verb|u8| to \verb|u64|, with
  \verb|wrap| where C goes round \\
\verb|char|, \verb|short|, \verb|int|, \verb|long| & \verb|int|, and a
  store into a smaller type is checked \\
\verb|size_t| & \verb|int|, where the value is never below zero;
  else kept \\
\verb|float|, \verb|double| & \verb|f32|, \verb|float| \\
\verb|char *| holding a text & \verb|str| \\
\verb|struct| & a class \\
an array & a list \\
\bottomrule
\end{tabular}
\end{center}

A signed sum that goes past the top is undefined in C, so no answer is
C's: Zeno stops with \verb|OverflowError|, and the report says where a
store into a smaller type is checked. \verb|/| and \verb|%| of whole
numbers that may be below zero go through \verb|c_divide| and
\verb|c_remainder|, which take the quotient toward zero; when both
sides are never below zero they are \verb|//| and \verb|%|.

A hash of 64 bits reads as it does in C:

\begin{zenofrag}
hash ^= u64(u8(letter_code(text, i)))
hash = u64.wrap(hash * FNV_PRIME)
\end{zenofrag}

\begin{note}
A \verb|size_t| is an \verb|int|: it counts bytes and items, and
stands beside \verb|int|s on nearly every line. It is never past
$2^{63}$; if it were, Zeno would stop with \verb|OverflowError| where C
goes on.
\end{note}

\section{Structs and pointers}

A struct is a value in C and a class is shared in Zeno. So the
converter writes a copy, \verb|a.copy()|, wherever C copies:
\verb|b = a;|, a struct put into an array, or given to a function that changes
it. Where nothing can tell, as for a struct a function only reads, it
writes none.

Pointers are of a few kinds, and the converter writes four of them:

\begin{center}
\begin{tabular}{lP{7cm}}
\toprule
A pointer that & Is written as \\
\midrule
is indexed, or is an array parameter & a list \\
points to one number written through & a list of one \\
points to a struct & the object \\
is \verb|const char *| & a text \\
\bottomrule
\end{tabular}
\end{center}

Everything else a pointer does is kept: a pointer variable, pointer
sums, \verb|NULL|, \verb|malloc|, a list made of \verb|next| pointers.
C does those well, and they have no Zeno that would read as Zeno.

\section{What is kept, and what crosses}

C can call only C, so a function that kept C calls is kept too. A
variable lives on one side, so a global that kept C uses is kept, and
so is Zeno that uses it. The Zeno calls the block with numbers, floats,
bools and texts that C does not change; a call that would give it a
struct or an array is kept. The block gets every \verb|#include|,
\verb|#define|, struct and enum of the file, so the kept C is compiled
as it was.

A line of the preprocessor that chooses (\verb|#if|, \verb|#ifdef|)
keeps the whole file, since it is two programs and the converter would
have to pick one. A \verb|main(int argc, char **argv)| is kept and
given the program's arguments, so \verb|zeno run tool.zeno a b| gives
it what \verb|./tool a b| would.

\section{What it does not promise}

What C leaves undefined: a signed sum past the top and an index past
the end make Zeno stop, and a variable C left without a value starts
at 0. The order C works out the arguments of one call in is left to
the compiler; Zeno's is left to right. \verb|long| is read as 64 bits,
as everywhere but Windows. The rest is in \verb|docs/convert-c.md|.

\section{How it is checked}

The programs of C in \verb|tools/convert| are each built with
\verb|cc| and run, converted and run with \verb|zeno run|, and what the
two print and the status they end with are compared. Among them are a
ring buffer, a frame parser and a state machine as they are written
for a board, and all of that is written as Zeno, with nothing kept.
As for Python, the converter reads what it wrote with Zeno's own
checker before it writes the file.
